\scanpage{151}%
\mspage{67}%
% archive sheet number: 150
ment générateur dans $X^\wedge$ i.e. que \unsure{tt} objet $\deleted{\unsure{E}}F$ de $X^\wedge$ est limite inductive $F=\underset{I}{\underrightarrow{\lim}}\,F_i$, les $F_i$ \unsure{ds|dans} $X$ (on prend $I=X_{/F}$), d'où

\[
\begin{aligned}
\deleted{\mathrm{Hom}(F\unsure{,}}\ \mathrm{Hom}(F,\alpha(\unsure{\zeta|\xi}))
&\simeq\underset{I}{\underleftarrow{\lim}}\ \underbrace{\mathrm{Hom}(F_i,\alpha(\zeta))}_{\mathrm{Hom}(\unsure{\rho|\varphi}(F_i),\zeta)}\\
&\simeq\mathrm{Hom}(\underrightarrow{\lim}\,\rho(F_i),\zeta),
\end{aligned}
\]

donc le \unsure{foncteur} $\zeta\mathrel{\unsure{\mapsto|\vdash}}\mathrm{Hom}(F,\alpha(\zeta))$ est représenté par $\rho(F)\overset{\mathrm{def}}{=}\underrightarrow{\lim}\,\rho(F_i)$.

Pour prouver $c\Longrightarrow a$, \deleted{il suffit de prouver que c) implique} \inserted{\unsure{Pour} l'existence} des \inserted{petites} $\underrightarrow{\lim}$\inserted{ dans $M\boxtimes N$, et l'existence d'une partie \unsure{strictmt} génératrice}, on invoque le lemme général

\noindent\begin{tabular}{@{}|p{.93\linewidth}@{}}
\uline{Lemme} 4.3.3. Soit $\alpha:\unsure{\mathcal E|E}\hookrightarrow\unsure{\mathcal F|F}$ un foncteur \unsure{pl} fidèle, admettant un adjoint à gauche \inserted{$\rho$}. \inserted{\textcircled{a}} Si les $\underrightarrow{\lim}$ d'un type donné $I$ existent dans $\mathcal F$, elles existent dans $\mathcal E$, et se calculent comme
\end{tabular}

\[
\underset{I\deleted{\unsure{(\mathcal E)}}}{\underrightarrow{\lim}}_{\!(\mathcal E)}\,\zeta_i
\simeq\rho\bigl(\underset{I}{\underrightarrow{\lim}}_{\!\mathcal F}\,\alpha(\zeta_i)\bigr)
\tag*{(4.3.3.1)}
\]

\inserted{\textcircled{b} Si $U$ est une partie \inserted{(petite)} str. génératrice de $\mathcal F$, $\rho(U)$ est une petite p. str. gén. dans $\mathcal E$.}

En effet, \unsure{si} $\zeta$ désigne le deuxième \unsure{mb}, et $y$ un objet variable de $\unsure{\zeta|\mathcal E}$, on aura

\[
\begin{aligned}
\mathrm{Hom}_{\mathcal E}(\zeta,y)&\overset{\mathrm{adj}}{\simeq}
\mathrm{Hom}_{\mathcal F}\bigl(\underset{I}{\unsure{\underline{\lim}|\underrightarrow{\lim}}}_{\!\mathcal F}\,\alpha(\zeta_i),\alpha(y)\bigr)\\
&\simeq\underset{I}{\underleftarrow{\lim}}\ \underbrace{\mathrm{Hom}_{\mathcal F}(\alpha(\zeta_i),\alpha(y))}_{\simeq\mathrm{Hom}_{\mathcal E}(\zeta_i,y)\ \text{\unsure{car} }\alpha\ \text{pl. fid.}}
\end{aligned}
\]

